\chapter{Introduction}

\section{Motivation}
The use of robots in production and industry has been growing steadily since the 1950's\cite{hockstein_history_2007}. With the rise of robots in industrial applications, we have also witnessed robots becoming a part of modern-day life. Many of us utilize robots and robotic solutions to help us with household tasks such as lawn mowing or cleaning our homes. The many advancements of robotics have also given birth to legged and other bio-inspired robots, making the possibilities even greater.

Legged robotic platforms have gained widespread popularity and are regarded as adaptable robots capable of traversing rough and challenging terrains with their human-like mobility. With this, we have seen remarkable achievements in the capabilities of robot agility. We now have dynamic legged systems that can run, jump, dance, and climb \ \cite{8913831}\cite{doi:10.1126/scirobotics.adi7566}. 
As robots become more advanced and competent, they are more often deployed beyond isolated environments, where they often interact directly with humans. In these environments, we see a wide variety of applications with a greater focus on human-robot interaction. Such applications include companion robots\cite{companion_1}\cite{companion_2}, assistive robots\cite{Assistive}, art \cite{sun2016cant}, and entertainment\cite{Disney}. With applications centered around human interaction, we are faced with new challenges in the scope of robotic design in addition to satisfying the already complex requirements for kinematics, locomotion, and stability. The evaluation of the expressive capabilities of the robot would now rely on the subjective perception in human trials.

When we think of modern robots, we often think of hard, mechanical, and cold structures or tools rather than companions capable of expressing intent and emotion. The interaction and communication of physical robots and humans is multidisciplinary and presents issues that have been traditionally studied in psychology. The design of a robot's behavior, appearance, and social ability is therefore highly challenging. When designing robots that interact with humans, we often assume that the robot should resemble humaneness as closely as possible. However, when humans are faced with such robots, we are more likely to feel discomfort as they fall into the uncanny valley for negative behavioral intention\cite{stafford2010improved}\cite{woods2004desig}\cite{mori_uncanny_2012}. In order to foster intuitive and meaningful interactions, robot design must account for both physical appearance and behavioral expressiveness. A great medium for displaying intent, emotion, and expressiveness is through body language and motion\cite{beck2012emotional}. Designing robot motion to be driven by emotion and intent rather than pure effectiveness and functionality, we may achieve more natural human-robot interactions.

\section{Goal of thesis}

In this thesis, the aim is to construct and design a biped robot agent with high articulation capable of displaying emotion and intention through expressive and dynamic motion. When designing the robot, we will not strive for the very best of mechanical design but instead focus on simplicity while still satisfying the goal of expressing emotion and intent. 

Construction and design will be performed using commonly available tools and technologies such as 3D printing and off-the-shelf hardware. By doing this, we allow for rapid prototyping and enable hobbyists, students, engineers, and other researchers to replicate and contribute further to the project with ease. In line with this, the entire project, including CAD files, source code, and documentation, will be released as open-source.

The goal of the work surrounding this thesis can be split into 2 independent goals:

The first goal is to design and construct a highly articulated bipedal robot platform capable of expressive motion. Domestic robots are designed purely to be simple and effective, leaving little room for expressive capabilities. Many open source robotic platforms lack the complexity and articulation needed to effectively communicate emotion and intent in a way that is convincing and pleasant without coming across as uncanny. The final robot should be highly capable, articulated, and robust while still being inviting and trustworthy in appearance.

The second goal is to develop a framework for generating and executing expressive motion. By developing a robust framework for expressive motion, we hope to contribute to the field of expressive robotics and enhance trust, safety, social acceptance, and general interaction for robots that work and interact closely with humans. The final result should be a framework that generates a diverse range of motions capable of clearly conveying emotion and intent. These expressive motions should be sufficiently distinct to be reliably recognized by independent observers during human trials.

The work of Grandia et al. heavily inspires the final robot design presented in this thesis in \cite{Disney}, where a biped robot character was designed and enhanced with expressive motion.

\section{Research Objectives}
The following research goals the work presented in this thesis 
objectives:

\begin{enumerate}
    \item Design and construct a highly articulated bipedal robot platform capable of producing a wide range of expressive motions using accessible and commonly available manufacturing methods and hardware components.

    \item Develop a software framework for generating and executing expressive motion on the platform, capable of producing a diverse range of motions that clearly communicate distinct emotions and intents.

    \item Evaluate the expressive capabilities of the platform through a human perception study, assessing the degree to which independent observers can reliably identify the intended emotion from the robot's motion alone.
\end{enumerate}

\section{Outline}
Outline comes here